\documentclass[11pt,a4paper]{article}

\usepackage[utf8]{inputenc}
\usepackage[T1]{fontenc}
\usepackage[french]{babel}
\usepackage{amsmath,amssymb,amsthm}
\usepackage[ruled,vlined,linesnumbered]{algorithm2e}
\usepackage[margin=2.5cm]{geometry}
\usepackage{booktabs}
\usepackage{enumitem}
\usepackage{multirow}
\usepackage[hidelinks]{hyperref}

\renewcommand{\algorithmcfname}{Algorithme}
\SetKwInput{KwIn}{Entrées}
\SetKwInput{KwOut}{Sorties}
\SetKwFor{For}{pour}{faire}{fin pour}
\SetKwFor{While}{tant que}{faire}{fin tant que}
\SetKwFor{ForEach}{pour chaque}{faire}{fin pour chaque}
\SetKwIF{If}{ElseIf}{Else}{si}{alors}{sinon si}{sinon}{fin si}
\SetKwBlock{Loop}{boucler}{fin boucle}
\SetKw{Return}{retourner}
\SetKw{KwTo}{à}
\SetKw{Break}{sortir}
\SetKwComment{Comment}{$\triangleright$\ }{}
\DontPrintSemicolon

\theoremstyle{plain}
\newtheorem{theorem}{Théorème}
\newtheorem{lemma}[theorem]{Lemme}
\newtheorem{proposition}[theorem]{Proposition}
\theoremstyle{remark}
\newtheorem{remark}[theorem]{Remarque}

\newcommand{\Zk}{Z_k}
\newcommand{\Nk}{N_k}
\newcommand{\Dk}{D_k}
\newcommand{\Ecal}{\mathcal{E}}
\newcommand{\Scal}{\mathcal{S}}
\newcommand{\Rcal}{\mathcal{R}}
\newcommand{\ILP}{\textsc{Pline}}
\newcommand{\qs}{q^{\star}}

\title{\bfseries Une matheuristique certifiée pour l'optimisation d'une
fonction fractionnaire sur l'ensemble efficace\\d'un programme linéaire
fractionnaire multi-objectif en nombres entiers}
\author{}
\date{}

\begin{document}
\maketitle
\vspace{-2\baselineskip}

\begin{abstract}
\noindent
On considère le problème (P) : maximiser une fonction linéaire fractionnaire
$f$ sur l'ensemble $\Ecal$ des solutions efficaces d'un programme linéaire
fractionnaire multi-objectif en nombres entiers (MOILFP). Les méthodes
exactes de la littérature résolvent ce problème par énumération implicite de
$\Ecal$ ; elles ne fournissent aucune garantie lorsqu'elles sont interrompues,
alors même que l'énumération de $\Ecal$ devient impraticable dès que le front
s'épaissit.

Nous proposons une \emph{matheuristique certifiée} : une recherche
heuristique dans l'espace des critères, dont chaque sous-problème est un
programme en nombres entiers résolu exactement et dont chaque incumbent est
certifié efficace, couplée à un mécanisme de certification qui convertit un
budget de calcul borné en une borne supérieure \emph{valide} sur l'optimum.
La méthode encadre donc $\qs$ à tout instant, ce qu'aucune méthode exacte du
domaine ne fait.

Trois résultats structurent la construction. Le test d'efficacité du
théorème~\ref{th:eff} est \emph{entièrement entier}, donc exact sans
tolérance numérique, et fournit gratuitement un certificat de dominance
réutilisable comme mouvement de recherche. La coupe de dominance du
théorème~\ref{th:cut} transpose au cas fractionnaire la coupe « $+1$ » du cas
linéaire entier, sans qu'aucun pas minimal ne soit à estimer. Le
théorème~\ref{th:bound} convertit toute borne supérieure du sous-problème de
Dinkelbach en borne sur $\qs$, avec un dénominateur restreint à la région où
la borne agit.

Sur un lot systématique de 90 instances, tous les résultats des deux méthodes
passent une vérification indépendante ; la matheuristique rend la meilleure
valeur sur \textbf{90 instances sur 90} et \emph{prouve} l'optimalité sur
\textbf{84}, contre 76 et 73 pour la méthode exacte de référence. Une
campagne contrôlée de 216 instances établit en outre que la taille et la
corrélation entre critères agissent sur des axes \emph{orthogonaux} : la
première gouverne l'échec de la méthode exacte, la seconde la fermeture de la
borne. Les résultats négatifs mesurés sont rapportés au même titre que les
positifs.
\end{abstract}

\tableofcontents

\bigskip

\section{Introduction}

\subsection{Le problème et sa difficulté propre}

Un programme linéaire fractionnaire multi-objectif en nombres entiers
s'écrit
\[
  \text{(MOILFP)}\qquad
  \max\ \Zk(x)=\frac{c_k^{\!\top}x+a_k}{d_k^{\!\top}x+b_k},\quad k=1,\dots,p
  \qquad\text{s.c.}\quad Ax\le b,\ x\in\mathbb{Z}^n_+ .
\]
Les fonctions fractionnaires sont des mesures de performance naturelles --
rendement, productivité, ratio coût/bénéfice -- ce qui explique leur usage en
planification de production et en analyse financière. Notons
$\Scal=\{x\in\mathbb{Z}^n_+ : Ax\le b\}$ le domaine réalisable et
$\Ecal\subseteq\Scal$ l'ensemble des solutions \emph{efficaces} de (MOILFP),
c'est-à-dire des $x$ pour lesquels il n'existe aucun $y\in\Scal$ vérifiant
$Z(y)\ge Z(x)$ et $Z(y)\ne Z(x)$.

Le problème traité ici est
\[
  \text{(P)}\qquad \max\ f(x)=\frac{c^{\!\top}x+a}{d^{\!\top}x+b}
  \qquad\text{s.c.}\quad x\in\Ecal ,
\]
soit l'optimisation d'une fonction d'utilité fractionnaire \emph{sur
l'ensemble efficace}. Cette formulation modélise la situation d'un décideur
qui, parmi les compromis admissibles, en choisit un selon un critère propre.

La difficulté est double, et il importe de la nommer précisément.
D'une part $\Ecal$ n'est pas décrit par un système d'inégalités : il est
défini par une condition d'optimalité, ce qui interdit d'écrire (P) comme un
programme mathématique standard. D'autre part $\Ecal$ n'est ni convexe ni
connexe, et son cardinal peut croître très vite : les mesures rapportées à la
section~\ref{sec:campagne} montrent un rapport $|\Ecal|/|\Scal|$ variant de
$0{,}001$ à $0{,}34$ sur des instances de même taille.

\subsection{Ce que la littérature fait, et ce qu'elle ne fait pas}

Les méthodes exactes existantes procèdent par énumération implicite :
elles balaient les points entiers du domaine, testent leur efficacité, et
resserrent progressivement la région de recherche par des coupes. Le travail
de référence pour le cas fractionnaire multi-objectif est celui de Zerdani et
Moulaï~\cite{zerdani2011}, qui optimise une fonction \emph{linéaire} sur
l'ensemble efficace entier d'un MOILFP ; Drici et
collaborateurs~\cite{drici2018} prolongent cette ligne.

Ces méthodes partagent une caractéristique structurante : elles sont de type
\emph{tout ou rien}. Interrompues avant terme, elles ne fournissent ni
solution garantie efficace, ni borne sur l'écart à l'optimum. Or nos mesures
(section~\ref{sec:echelle}) montrent qu'à partir de $n=12$ la méthode exacte
ne conclut plus que sur une fraction des instances dans un budget réaliste.
Le régime où elle échoue est précisément celui où le décideur aurait le plus
besoin d'une réponse.

\subsection{Contribution}

Nous construisons une matheuristique dont la propriété centrale est la
suivante : \emph{à tout instant, la méthode dispose d'un encadrement valide}
\[
  q_{\text{lb}}\ \le\ \qs\ \le\ q_{\text{ub}} ,
\]
où $q_{\text{lb}}=f(x_{\text{best}})$ avec $x_{\text{best}}$ \emph{certifié
efficace}, et où $q_{\text{ub}}$ découle du théorème~\ref{th:bound}. Les
contributions sont :

\begin{enumerate}[leftmargin=*,itemsep=2pt]
\item \textbf{Un test d'efficacité intégral} (th.~\ref{th:eff}), exact sans
      tolérance numérique, coûtant un seul programme en nombres entiers, et
      fournissant gratuitement un certificat de dominance.
\item \textbf{Une coupe de dominance fractionnaire exacte} (th.~\ref{th:cut}),
      transposition de la coupe « $+1$ » du cas linéaire entier, avec un
      big-M renforcé par relaxation continue.
\item \textbf{Un théorème de certification} (th.~\ref{th:bound}) qui convertit
      une borne supérieure du sous-problème de Dinkelbach en borne sur $\qs$,
      en restreignant le dénominateur à la région où la borne agit.
\item \textbf{Un protocole de validation à trois niveaux}, dont un niveau qui
      \emph{ne requiert aucune vérité terrain} et lève ainsi le plafond de
      taille imposé par l'énumération exhaustive.
\item \textbf{Une campagne contrôlée} qui établit l'orthogonalité de deux
      axes de difficulté, et un ensemble de \textbf{résultats négatifs}
      rapportés au même titre que les positifs.
\end{enumerate}

\section{Notations et hypothèses}

On note $\Nk(x)=c_k^{\!\top}x+a_k$ et $\Dk(x)=d_k^{\!\top}x+b_k$ les
numérateur et dénominateur du critère $k$, et $N(x)$, $D(x)$ ceux de la
fonction d'utilité $f$. On pose $\qs=\max_{x\in\Ecal}f(x)$. Un paramètre
rationnel $q$ est toujours écrit $q=P/Q$ sous forme irréductible avec $Q>0$.

\paragraph{Hypothèses.}
\begin{description}[leftmargin=1.4cm,style=nextline,itemsep=1pt]
\item[(A1)] $d_k\ge 0$ et $b_k\ge 1$ pour tout $k$, ainsi que pour $f$.
\item[(A2)] $A\ge 0$, à colonnes non nulles, et $b\ge 0$.
\end{description}

\noindent
(A1) entraîne $\Dk(x)\ge 1>0$ sur tout le domaine : les dénominateurs ne
s'annulent jamais et gardent un signe constant, ce qui est la condition de
validité des théorèmes~\ref{th:eff} et~\ref{th:cut}. (A2) entraîne que
$\Scal$ est non vide ($x=0$ est réalisable) et borné, avec les bornes
explicites
\[
  \overline{x}_j=\min_{i\,:\,A_{ij}>0}\left\lfloor b_i/A_{ij}\right\rfloor .
\]
Ces hypothèses ne sont pas des commodités techniques : ce sont elles qui
rendent la suite exacte. Toute instance extérieure doit être testée contre
elles avant traitement.

\paragraph{Convention d'intégralité.}
Toutes les données étant entières, les quantités $\theta$, $e_k$ et $F(q)$
introduites plus bas sont \emph{entières}. Les tests d'arrêt portant sur elles
sont donc exacts, sans $\varepsilon$ : c'est un point de mise en \oe uvre
essentiel, qu'une formulation directe en $v_k=\Nk/\Dk$ flottant perdrait.

\section{Briques exactes}

\subsection{Linéarisation de seuil}

\begin{theorem}[Linéarisation de seuil]\label{th:seuil}
Sous (A1), pour $v=N/D$ rationnel avec $D>0$ entier,
\[
  \Zk(x)\ \ge\ v
  \iff
  (D\,c_k-N\,d_k)^{\!\top}x\ \ge\ N b_k-D a_k .
\]
\end{theorem}

\begin{proof}
$\Zk(x)\ge v$ s'écrit $\Nk(x)\,D\ge N\,\Dk(x)$ après multiplication par
$D\,\Dk(x)>0$, quantité strictement positive sous (A1). En développant et en
regroupant, on obtient l'inégalité annoncée.
\end{proof}

\noindent
\emph{Portée.} Les contraintes de type $\varepsilon$-contrainte sont donc
\emph{entièrement linéaires et à coefficients entiers}. C'est ce qui permet
d'employer un solveur en nombres entiers standard à chaque appel, au lieu d'un
solveur fractionnaire dédié.

\subsection{Test d'efficacité intégral}

\begin{theorem}[Test d'efficacité]\label{th:eff}
Soit $\bar x\in\Scal$, $\bar N_k=\Nk(\bar x)$ et $\bar D_k=\Dk(\bar x)>0$.
Posons
\[
  \theta(\bar x)=\max\ \sum_{k=1}^{p}\Bigl[(\bar D_k c_k-\bar N_k d_k)^{\!\top}x
      +\bar D_k a_k-\bar N_k b_k\Bigr]
\]
sous les contraintes
$(\bar D_k c_k-\bar N_k d_k)^{\!\top}x\ge \bar N_k b_k-\bar D_k a_k$ pour tout
$k$, $Ax\le b$, $x\in\mathbb{Z}^n_+$. Alors $\theta(\bar x)$ est un entier
positif ou nul, et
\[
  \bar x\in\Ecal\iff\theta(\bar x)=0 .
\]
De plus, si $\theta(\bar x)>0$, tout optimum $x^\ast$ du programme domine
$\bar x$.
\end{theorem}

\begin{proof}
Le terme $k$ de l'objectif vaut, en développant,
\[
  \bar D_k\Nk(x)-\bar N_k\Dk(x)=\bar D_k\,\Dk(x)\,\bigl(\Zk(x)-\Zk(\bar x)\bigr).
\]
Comme $\bar D_k>0$ et $\Dk(x)>0$ sous (A1), ce terme est du signe de
$\Zk(x)-\Zk(\bar x)$. Les contraintes imposent que chacun de ces termes soit
positif ou nul, c'est-à-dire $Z(x)\ge Z(\bar x)$. Le point $\bar x$ étant
réalisable pour ce programme et y annulant l'objectif, on a
$\theta(\bar x)\ge 0$. Si $\theta(\bar x)=0$, aucun $x$ réalisable ne vérifie
$Z(x)\ge Z(\bar x)$ avec une inégalité stricte quelque part, donc $\bar x$
est efficace. Réciproquement, si $\theta(\bar x)>0$, un optimum $x^\ast$
vérifie $Z(x^\ast)\ge Z(\bar x)$ et $Z(x^\ast)\ne Z(\bar x)$ : il domine
$\bar x$, qui n'est donc pas efficace. L'intégralité de $\theta$ découle de
celle de toutes les données.
\end{proof}

\begin{remark}
Deux conséquences de mise en \oe uvre. D'abord, le test est \emph{exact} :
$\theta$ étant entier, le critère « $\theta=0$ » ne demande aucune tolérance.
Ensuite, en cas de non-efficacité, le programme fournit \emph{gratuitement}
un point dominant : cette information, qui serait perdue par un test binaire,
est exploitée à la fois comme mouvement de recherche locale et comme base de
coupe (théorème~\ref{th:cut}).
\end{remark}

\begin{remark}[Interruption]\label{rem:interruption}
Si le solveur est interrompu, deux cas seulement se présentent. Si
l'incumbent donne $\theta\ge 1$, alors $\bar x$ est \emph{prouvé} non
efficace et le certificat de dominance reste valide : un point dominant
suffit à conclure, il n'est pas nécessaire de connaître le maximum. Si
l'incumbent donne $\theta=0$, on ne peut rien conclure, car $\theta$ n'est
minoré que par $0$. Le verdict est alors \emph{indéterminé} et ne doit jamais
être lu comme « dominé » : ce test est le seul certificat d'efficacité de la
méthode, et un incumbent déclaré efficace à tort rendrait la borne inférieure
optimiste, c'est-à-dire la méthode inutilisable.
\end{remark}

\subsection{Dinkelbach à paramètre rationnel}

\begin{theorem}[Convergence finie]\label{th:dink}
Soit $X$ un ensemble fini non vide et
$F(q)=\max_{x\in X}\{N(x)-q\,D(x)\}$ avec $D>0$ sur $X$. Alors $F$ est le
supremum d'un nombre fini de fonctions affines de $q$ : elle est convexe,
linéaire par morceaux, strictement décroissante, et admet une racine unique
$q^{\ast}_X=\max_{x\in X}f(x)$. L'itération $q_{t+1}=f(x_t)$, où $x_t$
réalise $F(q_t)$, converge en un nombre \emph{fini} d'étapes.
\end{theorem}

\begin{proof}
Chaque $x\in X$ contribue l'affine $q\mapsto N(x)-qD(x)$, de pente
$-D(x)<0$ ; $F$ est donc convexe, linéaire par morceaux et strictement
décroissante, d'où l'unicité de la racine. Elle vaut $0$ exactement lorsque
$\max_X\{N(x)-qD(x)\}=0$, c'est-à-dire $q=\max_X f$. La finitude de $X$
borne le nombre de morceaux, donc le nombre d'itérations de Newton.
\end{proof}

\noindent
Le point important est que la convergence est \emph{finie} et non seulement
superlinéaire comme dans le cas continu : cela découle exactement de la
finitude de $\Ecal$. Avec $q=P/Q$ rationnel, le sous-problème s'écrit
\[
  \max\ (Q\,c-P\,d)^{\!\top}x+(Qa-Pb)
\]
à coefficients entiers, de valeur optimale entière : le critère d'arrêt
$F(q)\le 0$ est testé exactement.

\begin{remark}[Conséquence structurante]
Chaque itération revient à \emph{maximiser une fonction linéaire sur
l'ensemble efficace entier} -- exactement le problème traité
par~\cite{zerdani2011}. L'oracle interne dont nous avons besoin existe donc
déjà dans la littérature ; notre apport porte sur son emploi \emph{anytime}
et sur la certification.
\end{remark}

\section{Coupe de dominance et oracle}

\subsection{Coupe de dominance fractionnaire}

\begin{theorem}[Coupe de dominance]\label{th:cut}
Soit $\bar x\in\Scal$ \emph{dominé}. Posons
\[
  e_k(x)=\bar D_k\Nk(x)-\bar N_k\Dk(x)
        =\gamma_k^{\!\top}x+\delta_k
        =\bar D_k\,\Dk(x)\,\bigl(\Zk(x)-\Zk(\bar x)\bigr).
\]
Alors $e_k$ est à valeurs entières et
$\Zk(x)>\Zk(\bar x)\iff e_k(x)\ge 1$. De plus, la région
$\{x\in\Scal : Z(x)\le Z(\bar x)\}$ ne contient aucune solution efficace, et
peut donc être retirée sans perte.
\end{theorem}

\begin{proof}
L'identité et l'intégralité découlent du calcul déjà fait au
théorème~\ref{th:eff}. Comme $e_k$ est entière et de même signe que
$\Zk(x)-\Zk(\bar x)$, la stricte positivité équivaut à $e_k(x)\ge 1$ :
\emph{aucun pas minimal $\varepsilon$ n'est à estimer}, c'est la transposition
exacte au cas fractionnaire de la coupe « $+1$ » du cas linéaire entier.

Pour la seconde assertion, soit $y$ dominant $\bar x$, et soit $x$ tel que
$Z(x)\le Z(\bar x)$. Alors $Z(y)\ge Z(\bar x)\ge Z(x)$. Si $Z(y)=Z(x)$, alors
$Z(\bar x)$ est coincé entre les deux, d'où $Z(y)=Z(\bar x)$, ce qui
contredit le fait que $y$ domine $\bar x$. Donc $y$ domine $x$, et $x$ n'est
pas efficace.
\end{proof}

\noindent
C'est ce lemme qui rend la coupe sûre : \emph{on ne coupe jamais autour d'un
point efficace}, si bien que la question des solutions alternatives de même
vecteur critères ne se pose pas. La région est retirée par la disjonction
« il existe $k$ tel que $e_k(x)\ge 1$ », modélisée par $p$ binaires :
\[
  e_k(x)\ \ge\ 1-M_k(1-u_k),\qquad \sum_{k=1}^p u_k\ \ge\ 1,
  \qquad u\in\{0,1\}^p .
\]
En $\bar x$ on a $e_k(\bar x)=0$ pour tout $k$ : $\bar x$ est bien exclu.

\subsection{Renforcement du big-M}

Tout minorant de $\min_{x\in\Scal}e_k(x)$ fournit un $M_k$ valide. Le
minorant naïf, calculé sur la boîte $0\le x\le\overline{x}$, ignore
$Ax\le b$ et se révèle très lâche. Nous employons le minorant sur la
\emph{relaxation continue} de $\Scal$ : il est valide, puisque la relaxation
contient $\Scal$ ; il est plus fin, puisqu'elle est contenue dans la boîte ;
et, les coefficients étant entiers, on remonte au plafond entier du minorant
continu. Le coût est d'un programme linéaire par critère et par coupe,
comptabilisé séparément des programmes en nombres entiers.

\begin{table}[htbp]
\centering
\caption{Renforcement du big-M, 8 instances, 8 coupes chacune. $M_{\text{réel}}$
est calculé sur $\Scal$ \emph{énuméré} : la validité n'est pas postulée.}
\label{tab:bigm}
\begin{tabular}{lc}
\toprule
& médiane\\
\midrule
resserrement $M_{\text{boîte}}\to M_{\text{lp}}$ & $1{,}14\times$ (jusqu'à $1{,}47$)\\
écart résiduel $M_{\text{lp}}/M_{\text{réel}}$ & $1{,}00\times$ (au pire $1{,}05$)\\
validité $M_{\text{lp}}\ge M_{\text{réel}}$ & 8/8 instances\\
\bottomrule
\end{tabular}
\end{table}

\noindent
La seconde ligne du tableau~\ref{tab:bigm} est la plus informative : le
minorant continu \emph{atteint pratiquement le minimum entier}. Ce qui reste
à gagner sur le big-M lui-même est donc nul ou presque. C'est un résultat
utile bien que négatif : il ferme une piste avant qu'on y consacre du travail.

\subsection{Oracle anytime sur \texorpdfstring{$\Ecal$}{E}}
\label{sec:oracle}

L'oracle maximise une fonction linéaire $g$ sur $\Ecal$ en maintenant un
relâché $\Rcal\supseteq\Ecal$, initialisé à $\Scal$ et resserré par les
coupes du théorème~\ref{th:cut}.

\begin{algorithm}[H]
\caption{\textsc{MaxLinéaire}$(g,\Rcal,B)$ --- schéma}
\KwIn{objectif linéaire $g$ ; relâché $\Rcal$ ; budget $B$}
\KwOut{$(\textsf{statut},x_{\text{best}},\textsf{LB},\textsf{UB})$ avec
       $\textsf{LB}\le\max_{\Ecal}g\le\textsf{UB}$ à tout instant}
$\textsf{LB}\gets-\infty$ ;\quad $\textsf{UB}\gets+\infty$\;
\Loop{
  \lIf{budget épuisé}{\Return $(\textsf{limite},\cdot)$}
  $\bar x\gets\arg\max_{\Rcal} g$ ;\quad $\textsf{UB}\gets g(\bar x)$
    \Comment*[r]{un \ILP}
  \lIf{$\bar x$ efficace (th.~\ref{th:eff})}{\Return
    $(\textsf{optimal},\bar x,g(\bar x),\textsf{UB})$}
  $x_e\gets$ réparation de $\bar x$ par la chaîne de dominance\;
  $\textsf{LB}\gets\max(\textsf{LB},g(x_e))$\;
  \lIf{$\textsf{UB}-\textsf{LB}\le 0$}{\Return
    $(\textsf{optimal},x_{\text{best}},\textsf{LB},\textsf{UB})$}
  $\textsc{AjouterCoupe}(\Rcal,\bar x)$
    \Comment*[r]{licite : $\bar x$ est prouvé dominé}
}
\end{algorithm}

\begin{proposition}[Correction et terminaison]
L'invariant $\Ecal\subseteq\Rcal$ est préservé (th.~\ref{th:cut}). À l'arrêt
avec $\bar x$ efficace, $g(\bar x)=\max_{\Rcal}g\ge\max_{\Ecal}g$ et
$\bar x\in\Ecal$, donc $\bar x$ est optimal sur $\Ecal$. Chaque coupe retire
au moins $\bar x$ et $\Scal$ est fini : l'algorithme termine.
\end{proposition}

\noindent
\emph{Comportement anytime.} À chaque itération on dispose de
$\textsf{LB}\le\max_{\Ecal}g\le\textsf{UB}$. L'arrêt peut donc survenir à
tout instant avec un écart garanti, et survient souvent parce que
$\textsf{UB}$ rejoint $\textsf{LB}$, \emph{sans que la relaxation ait eu
besoin de produire un point efficace}. C'est le gain propre au schéma.

\subsection{Un point de sûreté}

L'assemblage naïf « Dinkelbach externe $+$ oracle interne » comporte un piège
que notre socle de vérité terrain a détecté. Lorsque l'oracle atteint sa
limite de temps, il renvoie $\textsf{LB}$, qui n'est qu'un \emph{minorant} de
$F(q)$. Traiter cette valeur comme exacte et conclure « $F(q)\le 0$ donc
optimal » produit des optima faux affichés comme prouvés. La règle correcte
est asymétrique :
\begin{itemize}[leftmargin=*,itemsep=1pt]
\item $\textsf{LB}>0$ reste valide quel que soit le statut, car
      $F(q)\ge\textsf{LB}>0$ : la racine n'est pas atteinte, on continue ;
\item $\textsf{LB}\le 0$ ne permet de conclure que si l'oracle a
      \emph{prouvé} l'optimalité ; sinon le statut « limite » est propagé.
\end{itemize}

\section{Certification : du budget à une garantie}

\subsection{Le théorème}

\begin{theorem}[Borne resserrée sur $\qs$]\label{th:bound}
Soit $q=P/Q$ atteinte par un point efficace (donc $q\le\qs$), soit
$\Rcal\supseteq\Ecal$ un relâché quelconque, et posons
\[
  U\ \ge\ F_{\Rcal}(q)=\max_{x\in\Rcal}\{Q\,N(x)-P\,D(x)\},
  \qquad
  D^{+}=\min\bigl\{D(x)\ :\ x\in\Rcal,\ Q\,N(x)-P\,D(x)\ge 0\bigr\}.
\]
Alors
\[
  \boxed{\ \qs\ \le\ q+\frac{U}{Q\,D^{+}}\ }
\]
et si $U=0$ alors $\qs=q$.
\end{theorem}

\begin{proof}
Soit $x\in\Ecal\subseteq\Rcal$. Si $Q N(x)-P D(x)<0$, alors $f(x)<q$ et $x$
ne contraint pas la borne. Sinon $x$ appartient à la région définissant
$D^{+}$, donc $D(x)\ge D^{+}$, et
\[
  f(x)-q=\frac{Q N(x)-P D(x)}{Q\,D(x)}\ \le\ \frac{U}{Q\,D^{+}} .
\]
En passant au maximum sur $\Ecal$ on obtient l'inégalité annoncée.
\end{proof}

\subsection{Trois gains sur la version naïve}

La version naïve du théorème prend $\Rcal=\Scal$ et
$D_{\min}=\min_{x\in\Scal}D(x)$. Le passage à la forme ci-dessus gagne sur
trois points, et il est instructif de les distinguer.

\begin{enumerate}[leftmargin=*,itemsep=3pt]
\item \textbf{$\Rcal$ au lieu de $\Scal$.} Les coupes retirent des zones sans
      aucune solution efficace : $U$ diminue.
\item \textbf{$D^{+}$ au lieu de $D_{\min}$.} On ne minimise $D$ que là où la
      borne agit, c'est-à-dire là où $f$ dépasse $q$. Le minimum portant sur
      un sous-ensemble, $D^{+}\ge D_{\min}$ mécaniquement. Mesures sur le lot :
      $D_{\min}\to D^{+}$ passe de $1$ à $12$, de $1$ à $13$, de $5$ à $10$,
      de $7$ à $24$.
\item \textbf{Recyclage des points dominés.} Les chaînes de réparation du
      théorème~\ref{th:eff} traversent des points \emph{dominés} -- les seuls
      sur lesquels le théorème~\ref{th:cut} autorise une coupe. Les jeter
      gaspillerait une information déjà payée. On les rejoue comme coupes,
      triés par valeur décroissante du substitut linéaire : ce sont eux qui
      tirent $U$ vers le haut, donc les couper resserre le plus.
\end{enumerate}

\begin{remark}
La région définissant $D^{+}$ n'est jamais vide : l'incumbent y appartient,
son résidu valant exactement $0$. Une infaisabilité signalerait donc un
défaut d'implémentation, non une preuve d'optimalité.
\end{remark}

\subsection{Effet mesuré}

\begin{table}[htbp]
\centering
\caption{Théorème naïf contre théorème~\ref{th:bound}, à budget identique,
8 instances du régime cible. Les deux colonnes sont produites par le même
code : seule la formule de la borne diffère.}
\label{tab:bound}
\begin{tabular}{lcc}
\toprule
& version naïve & théorème~\ref{th:bound}\\
\midrule
écart garanti médian & $31{,}4\ \%$ & $\mathbf{0{,}0\ \%}$\\
optimalité \emph{prouvée} & 0/8 & $\mathbf{5/8}$\\
validité $q_{\text{ub}}\ge\qs$ & 8/8 & 8/8\\
\bottomrule
\end{tabular}
\end{table}

\noindent
Le changement de nature compte davantage que le chiffre : sur cinq instances,
la matheuristique \emph{prouve} désormais l'optimalité, sur des instances que
la méthode exacte seule n'avait pas su résoudre dans le même budget. La
formulation « trouve sans pouvoir prouver » ne tient plus.

\section{La matheuristique}

\subsection{Architecture}

La méthode se compose de deux phases, dont la première ne peut jamais
compromettre la seconde.

\paragraph{Phase 1 -- recherche.} Une recherche à voisinage variable dans
l'\emph{espace des critères}. Chaque voisinage est défini par des contraintes
linéaires issues des théorèmes~\ref{th:seuil} et~\ref{th:cut}, et le
sous-problème correspondant est résolu \emph{exactement}. Chaque point obtenu
est ensuite certifié efficace par le théorème~\ref{th:eff}. Le choix des
régions est heuristique ; la résolution dans chaque région est exacte. C'est
le sens précis du mot « matheuristique » ici.

\paragraph{Phase 2 -- certification.} L'oracle de la
section~\ref{sec:oracle} est lancé avec un budget borné au point $q$ obtenu ;
sa borne supérieure $U$ est convertie en borne sur $\qs$ par le
théorème~\ref{th:bound}.

\subsection{Les trois mouvements, et une erreur de conception}

\paragraph{Une erreur instructive.} Une première version prenait pour
voisinage
\[
  V_k(x^r)=\{x\ :\ e_k(x)\ge 1,\ e_j(x)\ge 0\ \forall j\ne k\},
\]
soit « mieux sur $k$ sans rien perdre ailleurs ». Or c'est \emph{exactement}
l'ensemble des points qui dominent $x^r$ : il est \textbf{vide} dès que $x^r$
est efficace. Le voisinage était donc vide par construction, ce qui fut
vérifié sur instance, et la recherche ne progressait pas. Un arbitrage sur
les autres critères est indispensable.

\medskip
\begin{algorithm}[H]
\caption{Les trois mouvements}
\SetKwProg{Mv}{Mouvement}{ :}{}
\Mv{$A$ --- $\varepsilon$ partiel}{
  $J\gets$ sous-ensemble \textbf{strict} aléatoire de
    $\{1,\dots,p\}\setminus\{k\}$\;
  \Return $\arg\max\{w^{\!\top}x+w_0 : x\in\Scal,\ e_k(x)\ge 1,\
    e_j(x)\ge 0\ \forall j\in J\}$\;
}
\Mv{$B$ --- plancher absolu}{
  tirer $\varepsilon_j$ entre nadir$_j$ et idéal$_j$\;
  \Return $\arg\max\{w^{\!\top}x+w_0 : x\in\Scal,\ Z_j(x)\ge\varepsilon_j\}$
    \Comment*[r]{seuils linéarisés (th.~\ref{th:seuil})}
}
\Mv{$C$ --- LNS \emph{fix-and-optimize}}{
  fige une fraction des variables à leur valeur dans l'incumbent\;
  \Return l'optimum exact du sous-problème restant\;
}
\end{algorithm}

\medskip\noindent
La taille de $J$ règle l'intensification : $J$ vide donne le mouvement le
plus explorateur, $J$ plein redonne le voisinage vide ci-dessus.

\subsection{Redémarrages guidés par l'archive}

Une version antérieure arrêtait la recherche au premier tour sans
amélioration. Elle repart désormais des points de l'archive \emph{les moins
exploités} comme bases, et n'abandonne qu'après plusieurs tours stériles
consécutifs. L'archive ne contient que des points efficaces certifiés : ce
sont des bases légitimes et déjà payées, là où un redémarrage aléatoire
jetterait le travail accompli.

L'effet porte sur la \emph{dispersion} plus que sur la médiane : sur
l'instance la plus dure du lot ($|\Ecal|=218$), l'écart réel maximal sur
10 graines passe de $49{,}0\ \%$ à $0{,}0\ \%$.

\subsection{Allocation du budget et plafond de coupes}

Deux réglages qui étaient figés sont désormais arbitrés par le budget.
Le budget de recherche non consommé -- la recherche s'arrête sur stagnation --
est \emph{reversé} à la certification, qui en manquait. Et le plafond de
coupes recyclées, autrefois fixé à 40, est piloté par le budget : les coupes
sont ajoutées par lots, un lot de plus n'étant posé que si la borne n'a pas
fermé et qu'il reste du temps.

\begin{algorithm}[H]
\caption{\textsc{Matheuristique}$(B_{\text{rech}},B_{\text{cert}},s_{\max})$}
\KwIn{budgets de recherche et de certification ; seuil de stagnation}
\KwOut{$(q_{\text{lb}},x_{\text{best}},q_{\text{ub}})$ avec
       $q_{\text{lb}}\le\qs\le q_{\text{ub}}$}
$x_0\gets$ réparation de $0$ ;\quad $\mathcal{A}\gets\{x_0\}$ ;\quad
$q\gets f(x_0)$ ;\quad $s\gets 0$\;
\While{$\textsc{écoulé}()<B_{\text{rech}}$}{
  $(w,w_0)\gets$ substitut de Dinkelbach en $q$ ;\quad
  $\textsf{amélioré}\gets\textbf{faux}$\;
  $\mathcal{B}\gets\textsc{ChoisirBases}(\mathcal{A},\text{usage},\ s>0)$\;
  \ForEach{$x^r\in\mathcal{B}$, $k$ dans une permutation de $\{1..p\}$}{
    $y\gets\textsc{Mouvement}_{A|B|C}(x^r,k,w,w_0)$\;
    $y\gets$ réparation certifiée de $y$ (th.~\ref{th:eff})\;
    \lIf{$y=\bot$}{\textbf{continuer} \Comment*[f]{non certifié : rien n'entre}}
    $\mathcal{A}\gets\mathcal{A}\cup\{y\}$\;
    \lIf{$f(y)>q$}{$q\gets f(y)$ ;\ $x_{\text{best}}\gets y$ ;\
      $\textsf{amélioré}\gets\textbf{vrai}$}
  }
  $s\gets 0$ \textbf{si} $\textsf{amélioré}$ \textbf{sinon} $s+1$\;
  \lIf{$s\ge s_{\max}$}{\Break}
}
$B\gets B_{\text{cert}}+\max(0,\ B_{\text{rech}}-\textsc{écoulé}())$
  \Comment*[r]{réallocation}
$q_{\text{ub}}\gets\textsc{Certifier}(q,\mathcal{A},B)$ (th.~\ref{th:bound})\;
\Return $(q,x_{\text{best}},q_{\text{ub}})$\;
\end{algorithm}

\medskip\noindent
\emph{Propriété fondamentale.} L'incumbent est \emph{toujours} un point
efficace certifié : la borne inférieure n'est jamais optimiste, même si la
recherche est interrompue. C'est la seule propriété que la méthode n'a pas le
droit de perdre, et tout le reste du dispositif est organisé pour la
préserver -- en particulier la réparation rend $\bot$ plutôt qu'un point non
certifié (remarque~\ref{rem:interruption}).

\section{Protocole de validation}
\label{sec:validation}

Aucune brique n'est réputée correcte tant qu'elle n'a pas été confrontée à
une référence indépendante. Le protocole comporte trois niveaux, adaptés à ce
qui est vérifiable à chaque taille.

\subsection{Niveau 1 --- contre la vérité terrain ($n\le 8$)}

L'énumération exhaustive de $\Scal$, suivie d'un filtrage de Pareto en
arithmétique rationnelle exacte, fournit $\Ecal$ et $\qs$. On vérifie :
$\theta(x)=0\iff x\in\Ecal$ pour tout $x$ testé, et que le dominateur rendu
domine réellement ; Dinkelbach retrouve $\max_{\Scal}f$ ; l'oracle retrouve
$\max_{\Ecal}g$ pour plusieurs $g$ ; l'hybride retrouve $\qs$ ; l'archive ne
contient aucun faux positif ; et le big-M majore le minimum \emph{réel} de
$e_k$ sur $\Scal$ énuméré.

\subsection{Niveau 2 --- sans vérité terrain (jusqu'à $n=40$)}
\label{sec:w}

L'énumération plafonne vers $n=8$. Au-delà, les garanties se contrôlent
néanmoins, car elles sont \emph{auto-certifiantes}.

\begin{table}[htbp]
\centering
\caption{Contrôles ne requérant pas $\qs$.}
\begin{tabular}{clp{9.6cm}}
\toprule
& propriété & comment elle se vérifie sans $\qs$\\
\midrule
$W_1$ & $q_{\text{lb}}\le\qs$ & $x_{\text{best}}$ passe le test d'efficacité\\
$W_2$ & archive saine & même test sur un échantillon\\
$W_3$ & encadrement cohérent & $q_{\text{lb}}\le q_{\text{ub}}$\\
$W_4$ & coupes sûres & tout point efficace certifié satisfait la disjonction
        de chaque coupe posée : c'est $\Ecal\subseteq\Rcal$ restreint à la
        partie \emph{connue} de $\Ecal$, en arithmétique entière pure\\
$W_5$ & encadrements concordants & plusieurs exécutions indépendantes
        encadrent le même $\qs$, donc
        $\max_r q_{\text{lb}}^r\le\min_r q_{\text{ub}}^r$ ; une violation
        \emph{prouve} un défaut\\
$W_6$ & témoin indépendant & $\max_{\Scal}f$ est une borne supérieure valide
        et exacte, obtenue par un seul Dinkelbach, sans énumération\\
\bottomrule
\end{tabular}
\end{table}

\noindent
$W_4$ remplace le contrôle de sûreté des coupes que l'énumération assurait ;
$W_5$ transforme la simple répétition d'exécutions en test de cohérence
rigoureux. Le protocole établit la \textbf{validité} des bornes, non leur
finesse : une borne correcte mais inutile passe $W_1$--$W_6$ sans broncher.

\subsection{Niveau 3 --- comparaison entre méthodes}

Comparer suppose un lot d'instances identique, une interface commune, et
surtout une \emph{vérification indépendante} de ce que chaque méthode
renvoie : on ne compare pas des nombres auto-déclarés. Tout résultat, y
compris les nôtres, est recontrôlé : $x$ réalisable ($C_1$), $x$ efficace par
le théorème~\ref{th:eff} ($C_2$), $q=f(x)$ ($C_3$), et contre la vérité
terrain $q\le\qs$ avec égalité dès que le statut annonce l'optimalité
($C_4$). Le vérificateur est lui-même testé par une méthode délibérément
fautive, rejetée par $C_2$ et $C_4$ : un contrôle qu'on n'a jamais vu échouer
ne prouve rien.

\section{Campagne de difficulté}
\label{sec:campagne}

\subsection{Plan d'expérience}

Grille : $n\in\{5,6,7,8\}$, $p\in\{2,3,4\}$,
$\text{corr}\in\{0;0{,}25;0{,}5;0{,}75;0{,}9;1\}$, 3 graines, soit
\textbf{216 instances}, vérité terrain par énumération, limite de 12~s par
résolution.

Le paramètre $\text{corr}$ pilote la similarité entre critères, donc la
finesse de $\Ecal$, \textbf{à domaine $\Scal$ rigoureusement identique} : le
générateur construit $A$ et $b$ indépendamment de $\text{corr}$, ainsi que la
fonction d'utilité. On \emph{manipule} donc la cause supposée au lieu de
l'observer : le plan est causal et non corrélationnel.

\paragraph{Validation.} \textbf{0 divergence} avec la vérité terrain sur les
199 instances résolues à l'optimum, \textbf{0 faux positif} dans les
archives, 17 arrêts sur limite honnêtement rapportés.

\subsection{Traitement de la censure}

Les 17 instances arrêtées ont un coût \emph{censuré} : leur nombre d'appels
est un minorant, pas une mesure. Les inclure dans une moyenne revient à
remplacer le coût des instances les plus dures par une valeur tronquée, et
\emph{les strates qui échouent le plus paraissent alors les plus faciles}.
C'est un biais de survie. On sépare donc deux questions.

\subsection{Q1 --- coût parmi les instances résolues}

\begin{table}[htbp]
\centering
\caption{Corrélations de Spearman avec le nombre d'appels au solveur,
censurées exclues. \emph{***} : $p<10^{-3}$, \emph{*} : $p<5\cdot 10^{-2}$.}
\begin{tabular}{lr}
\toprule
prédicteur & Spearman\\
\midrule
$\text{relax\_gap}=(\max_{\Scal}f-\qs)/|\max_{\Scal}f|$ & $+0{,}607$ ***\\
$\theta_{\Scal}$ (un seul \ILP) & $+0{,}409$ ***\\
$\text{depth}_{\Scal}$ (quelques \ILP) & $+0{,}373$ ***\\
rapport $|\Ecal|/|\Scal|$ & $\mathbf{-0{,}352}$ ***\\
$|\Ecal|$ absolu & $\mathbf{-0{,}306}$ ***\\
$p$ & $-0{,}273$ ***\\
$n$ & $+0{,}268$ ***\\
$|\Scal|$ & $+0{,}173$ *\\
\bottomrule
\end{tabular}
\end{table}

\noindent
Le coût monte quand $\Ecal$ est \emph{mince} : la relaxation est alors loin de
$\Ecal$ et il faut beaucoup de coupes pour l'y ramener. Notons que
$\text{relax\_gap}$ fait intervenir $\qs$, donc la réponse cherchée : c'est
un indicateur de \emph{mécanisme}, jamais un prédicteur utilisable avant
résolution.

\subsection{Q2 --- échec total : le sens inverse}

\begin{table}[htbp]
\centering
\caption{Censurées contre résolues (Mann--Whitney bilatéral).}
\begin{tabular}{lrrr}
\toprule
variable & médiane censurée & médiane résolue & $p$\\
\midrule
$|\Ecal|$ & $\mathbf{26}$ & $\mathbf{4}$ & $1{,}2\cdot 10^{-3}$ **\\
$|\Ecal|/|\Scal|$ & $0{,}0139$ & $0{,}0017$ & $7{,}7\cdot 10^{-3}$ **\\
$n$ & $7$ & $6$ & $3{,}2\cdot 10^{-2}$ *\\
$|\Scal|$ & $2387$ & $1931$ & $5{,}3\cdot 10^{-2}$ (n.s.)\\
$\theta_{\Scal}$ & $3229$ & $2819$ & $0{,}42$ (n.s.)\\
$\text{depth}_{\Scal}$ & $1$ & $1$ & $0{,}80$ (n.s.)\\
\bottomrule
\end{tabular}
\end{table}

\noindent
\textbf{Deux régimes distincts, de sens opposés.} Un $\Ecal$ mince rend
chaque résolution modérément plus chère ; un $\Ecal$ \emph{gros} rend la
résolution impossible dans le temps imparti, car l'oracle doit couper un à un
un grand nombre de points non dominés et dérive vers un comportement
d'énumération.

Conséquence lourde : $\theta_{\Scal}$ et $\text{depth}_{\Scal}$, les deux
prédicteurs bon marché, corrèlent avec le coût des instances résolues mais
\textbf{n'ont aucun pouvoir de prédiction sur l'échec} ($p=0{,}42$ et
$p=0{,}80$). Ils détectent le facile, pas le catastrophique. C'est un
résultat négatif, et il est utile : il délimite ce qu'il reste à trouver.

\subsection{Q3 et Q4 --- corrélation et taille}

\begin{table}[htbp]
\centering
\begin{minipage}{0.48\textwidth}
\centering
\caption{Q3 : expérience contrôlée sur corr.}
\begin{tabular}{cccc}
\toprule
corr & $|\Ecal|$ méd. & timeouts & \ILP\ méd.\\
\midrule
$0{,}00$ & $44{,}5$ & \textbf{6} & $29{,}5$\\
$0{,}25$ & $47{,}0$ & \textbf{4} & $34{,}5$\\
$0{,}50$ & $18{,}5$ & \textbf{6} & $59{,}5$\\
$0{,}75$ & $2{,}0$ & 1 & $83{,}0$\\
$0{,}90$ & $1{,}0$ & 0 & $65{,}5$\\
$1{,}00$ & $1{,}0$ & 0 & $53{,}0$\\
\bottomrule
\end{tabular}
\end{minipage}
\hfill
\begin{minipage}{0.48\textwidth}
\centering
\caption{Q4 : la taille prédit-elle ?}
\begin{tabular}{crrrr}
\toprule
$n$ & méd. & min & max & \% timeout\\
\midrule
5 & 41 & 5 & 214 & $1{,}9$\\
6 & 44 & 7 & 236 & $5{,}6$\\
7 & 57 & 5 & \textbf{359} & $13{,}0$\\
8 & 92 & 10 & 215 & $11{,}1$\\
\bottomrule
\end{tabular}
\end{minipage}
\end{table}

\noindent
Des critères corrélés font \emph{s'effondrer} $\Ecal$ (médiane $44\to 1$) et
le timeout disparaît. Lue sur la seule colonne des coûts, la corrélation
semble au contraire rendre les instances plus chères ($29{,}5\to 83{,}0$), ce
qui n'est que l'effet du biais de survie. À $n$ fixé, le coût couvre deux
ordres de grandeur : rapporter les performances en fonction de la seule
taille ne permet aucune prédiction utile.

\paragraph{Confirmations.} Dinkelbach converge en 1 à 4 itérations externes
(médiane 1, moyenne $1{,}56$) sur les 216 instances : la convergence finie du
théorème~\ref{th:dink} est rapide en pratique. L'archive couvre $58{,}6\ \%$
de $\Ecal$ en médiane, sans aucun faux positif.

\section{Résultats}
\label{sec:resultats}

\subsection{Comparaison sur un lot systématique}

Le lot \texttt{lot\_v1} comprend 90 instances en deux strates : 72 à
$n\le 8$ avec vérité terrain, 18 à $n$ de 10 à 20 sans. Il s'agit d'une
grille systématique $n\times p\times\text{corr}\times$ graine : contrairement
aux lots choisis \emph{a posteriori} parce qu'une méthode y échouait,
l'argument n'est pas circulaire.

\begin{table}[htbp]
\centering
\caption{Lot figé, 90 instances, budget identique de 12~s par méthode.}
\label{tab:lot}
\setlength{\tabcolsep}{5pt}
\begin{tabular}{lccccc}
\toprule
& optimalité & meilleure & \ILP & $t$ & vérifications\\
& \textbf{prouvée} & valeur & méd. & méd. & KO\\
\midrule
méthode exacte & 73/90 & 76/90 & 92 & $1{,}5$~s & \textbf{0}\\
\textbf{matheuristique} & $\mathbf{84/90}$ & $\mathbf{90/90}$ & 124
& $\mathbf{0{,}7}$~s & \textbf{0}\\
\bottomrule
\end{tabular}
\end{table}

\begin{table}[htbp]
\centering
\caption{Détail par strate.}
\begin{tabular}{llcc}
\toprule
strate & & exacte & matheuristique\\
\midrule
\multirow{2}{*}{$n\le 8$ (72 inst., vérité terrain)}
 & optimalité prouvée & 65/72 & $\mathbf{70/72}$\\
 & meilleure valeur & 66/72 & $\mathbf{72/72}$\\
\midrule
\multirow{2}{*}{$n=10\dots20$ (18 inst., sans)}
 & optimalité prouvée & 8/18 & $\mathbf{14/18}$\\
 & meilleure valeur & 10/18 & $\mathbf{18/18}$\\
\bottomrule
\end{tabular}
\end{table}

\noindent
Trois lectures, dans l'ordre d'importance.
\begin{enumerate}[leftmargin=*,itemsep=2pt]
\item \textbf{Les 180 résultats passent $C_1$--$C_4$.} C'est ce qui donne son
      sens au reste : aucune valeur du tableau n'est auto-déclarée.
\item \textbf{La matheuristique rend la meilleure valeur sur les 90
      instances}, et pas seulement dans le régime où l'exact échoue.
\item \textbf{Elle prouve l'optimalité plus souvent que la méthode exacte}
      (84 contre 73). Ce résultat contre-intuitif tient au
      théorème~\ref{th:bound} : elle certifie sans avoir à fermer le
      sous-problème, là où la méthode exacte doit le résoudre.
\end{enumerate}
Le coût médian en appels est plus élevé (124 contre 92) pour un temps médian
plus faible ($0{,}7$~s contre $1{,}5$~s) : les sous-problèmes sont plus
nombreux mais individuellement plus petits. C'est pourquoi les deux unités
sont rapportées.

\subsection{Stabilité entre graines}

La recherche étant stochastique, un tableau à une seule graine illustre un
comportement mais ne le mesure pas. Sur 8 instances $\times$ 10 graines à
budget identique : écart garanti médian $\mathbf{0{,}00\ \%}$, optimum
atteint sur \textbf{78 exécutions sur 80}, optimalité prouvée sur
\textbf{57/80}, et l'encadrement $q_{\text{lb}}\le\qs\le q_{\text{ub}}$
vérifié sur \textbf{80/80}.

\subsection{Passage à l'échelle : deux axes orthogonaux}
\label{sec:echelle}

\begin{table}[htbp]
\centering
\caption{Écart garanti médian, grille $n\times\text{corr}$, 30~s par méthode.}
\label{tab:echelle}
\begin{tabular}{lccc}
\toprule
$n$ & $\text{corr}=0{,}00$ & $\text{corr}=0{,}50$ & $\text{corr}=0{,}90$\\
\midrule
10 & $0{,}0\ \%$ & $0{,}0\ \%$ & $0{,}0\ \%$\\
15 & $37{,}9\ \%$ & $0{,}0\ \%$ & $0{,}0\ \%$\\
20 & $61{,}0\ \%$ & $10{,}8\ \%$ & $0{,}0\ \%$\\
25 & $44{,}6\ \%$ & $4{,}6\ \%$ & $0{,}0\ \%$\\
30 & $82{,}4\ \%$ & $24{,}8\ \%$ & $42{,}6\ \%$\\
40 & $92{,}8\ \%$ & $70{,}8\ \%$ & $19{,}6\ \%$\\
\midrule
\textbf{toutes tailles} & $\mathbf{71{,}6\ \%}$ & $\mathbf{4{,}6\ \%}$
& $\mathbf{0{,}0\ \%}$\\
\bottomrule
\end{tabular}
\end{table}

\noindent
La méthode exacte, elle, prouve l'optimalité sur 3 instances sur 12 à
\emph{chacun} des trois niveaux de corr, avec la même décroissance 2/2, 1/2,
puis 0/2 partout dès $n=20$.

\begin{proposition}[Dissociation]
La taille gouverne l'échec de la méthode exacte, et la corrélation n'y change
rien. La corrélation gouverne la fermeture de la borne, et beaucoup plus que
la taille. Les deux axes sont donc \emph{orthogonaux}, et il faut deux
leviers distincts.
\end{proposition}

\noindent
L'effet protecteur de corr observé à petite taille ($n$ de 5 à 8, où les
17 échecs étaient tous à $\text{corr}\le 0{,}75$) \emph{disparaît} une fois
$n$ assez grand : au-delà, la taille suffit à faire échouer la méthode exacte
quelle que soit la finesse de $\Ecal$.

Une vérification directe le confirme. À $n=20$ et $n=25$, la méthode exacte
lancée 400~s ne conclut pas et rend $2{,}711$ et $3{,}625$ ; la
matheuristique, en 18~s, \emph{certifie} $3{,}732$ et $5{,}738$. La recherche
n'est pas le goulot : la borne l'est, dans le régime à $\Ecal$ épais.

\begin{remark}[Réserve]
Le tableau~\ref{tab:echelle} n'est pas monotone ($n=30$, $\text{corr}=0{,}9$ :
$42{,}6\ \%$ contre $19{,}6\ \%$ à $n=40$), et deux instances par cellule ne
permettent pas de trancher entre effet et bruit. Une validation prenant le
meilleur de trois exécutions donne $0{,}0\ \%$ sur cette même cellule : la
variance entre graines est au moins aussi grande que l'effet de la taille.
Conclusion à ne pas dépasser : corr domine, $n$ compte moins qu'il n'y paraît.
\end{remark}

\section{Résultats négatifs}

Ils sont rapportés au même titre que les positifs, parce qu'ils délimitent ce
qu'il reste à chercher.

\paragraph{Le nombre de coupes sature, puis nuit.} Une première version du
plafond adaptatif, à petits lots et nombreux tours, a \emph{dégradé} l'écart
garanti médian ($2{,}94\ \%\to 6{,}17\ \%$). Diagnostic par les journaux : le
modèle étant conservé d'un tour à l'autre, les coupes s'accumulent et
atteignaient environ 90 coupes, soit $p$ binaires chacune ; chaque
sous-problème devenait trop lourd pour que l'oracle referme quoi que ce soit,
alors même qu'il disposait de deux fois plus de temps.

\begin{table}[htbp]
\centering
\caption{Réglage du plafond de coupes, 4 instances $\times$ 4 graines.}
\begin{tabular}{lcc}
\toprule
schéma & écart garanti médian & optimalité prouvée\\
\midrule
lots de 10, 6 tours & $9{,}91\ \%$ & 4/16\\
lots de 40, 2 tours & $\mathbf{0{,}00\ \%}$ & 11/16\\
lot unique de 40 & $0{,}00\ \%$ & 10/16\\
lot unique de 80 & $0{,}00\ \%$ & 12/16\\
\bottomrule
\end{tabular}
\end{table}

\noindent
La leçon n'est pas le chiffre mais le sens : \textbf{plus de coupes n'est pas
mieux}. Tant que chaque coupe coûte $p$ binaires, en ajouter davantage se
paie. Désagréger la disjonction conditionne donc tout gain supplémentaire.

\paragraph{Le big-M est une piste fermée.} Le tableau~\ref{tab:bigm} le
montre : le minorant continu atteint pratiquement le minimum entier.

\paragraph{La limite de temps par sous-problème n'achète rien de mesurable.}
Le passage d'une limite de temps réelle au solveur -- correction pourtant
légitime, puisqu'un seul appel pouvait en principe absorber tout le budget --
n'a produit \emph{aucun} effet mesurable : $0{,}00\ \%$ d'écart médian,
78/80 optima, 57/80 preuves, avant comme après. Le dépassement réel
plafonnait à $+1{,}3$~s. C'est une \emph{assurance}, pas un résultat, et la
rapporter comme un gain serait trompeur.

\paragraph{Une fausse piste, et sa leçon de méthode.} En mesurant $f$ avant
et après la chaîne de réparation, on observe une perte de $74$ à $82\ \%$ :
le mouvement trouve $f=31{,}75$, la réparation rend $f=1{,}77$. Le test
d'efficacité maximise $\theta$, donc choisit le dominateur qui maximise
l'amélioration des \emph{critères}, en ignorant $f$. Deux réparations guidées
furent implémentées ; aucune ne fait mieux à $n=30$, et celle qui maximise
$f$ exactement coûte trente fois plus d'appels. La comparaison avec la
méthode exacte a expliqué pourquoi : les points à $f$ élevé sont
\emph{dominés}, donc hors de $\Ecal$, et leurs voisins efficaces ont
réellement un $f$ faible. \textbf{Une mesure locale spectaculaire ne vaut pas
diagnostic tant qu'un témoin indépendant ne l'a pas confirmée.}

\section{Limites}

Elles sont énoncées ici plutôt que laissées à découvrir.

\begin{enumerate}[leftmargin=*,itemsep=3pt]
\item \textbf{Aucun témoin issu de la littérature.} Les méthodes
      de~\cite{zerdani2011} et~\cite{drici2018} ne sont pas réimplémentées :
      les articles n'ont pu être obtenus dans notre environnement de
      développement, et réimplémenter de mémoire un algorithme non lu
      produirait un homme de paille, que l'on battrait par construction. La
      seule comparaison disponible reste donc \emph{interne}. C'est la limite
      dominante pour la publication, et le banc décrit au niveau~3 du
      protocole est prêt à accueillir ces témoins sans modification.
\item \textbf{Calibre modéré.} $n\le 20$ dans le lot comparatif, $n\le 40$
      dans la validation. La vérité terrain par énumération plafonne vers
      $n=8$, ce que le niveau~2 du protocole contourne mais ne supprime pas.
\item \textbf{Une seule famille d'instances.} Toutes proviennent d'un même
      générateur aléatoire sous (A1)--(A2). La transférabilité à des
      instances structurées ou réelles n'est pas établie.
\item \textbf{$p$ fixé à 3 à grande taille}, alors même que le coût des
      coupes croît en $p$ ; et deux instances par cellule dans l'étude
      d'échelle, quand la variance entre graines est du même ordre que
      l'effet mesuré.
\item \textbf{Un seul solveur.} Le comptage des appels atténue la dépendance
      sans la supprimer.
\end{enumerate}

\section{Conclusion et perspectives}

Nous avons construit une matheuristique qui, sur le problème d'optimisation
d'une fonction fractionnaire sur l'ensemble efficace d'un MOILFP, fournit à
tout instant un encadrement valide de l'optimum. La borne inférieure repose
sur un incumbent \emph{certifié efficace} par un test intégral exact ; la
borne supérieure découle du théorème~\ref{th:bound}, qui convertit un budget
borné en garantie. Sur un lot systématique de 90 instances, tous les
résultats passent une vérification indépendante, et la méthode rend la
meilleure valeur partout tout en prouvant l'optimalité plus souvent que la
méthode exacte de référence.

Deux directions se dégagent, dans cet ordre.

\paragraph{Resserrer la borne dans le régime à $\Ecal$ épais.} C'est le
verrou, établi de deux côtés indépendants : à petite taille sur une instance
où prouver $F(\qs)\le 0$ demande 433~s, 177 coupes et 533 appels, soit
43~fois le budget de certification ; et à grande taille par le
tableau~\ref{tab:echelle}. Le préalable est de \textbf{désagréger la coupe de
dominance} : tant que chaque coupe coûte $p$ binaires, aucune stratégie de
coupes ne fera mieux.

\paragraph{Établir la comparaison externe.} Réimplémenter~\cite{zerdani2011}
et~\cite{drici2018} sur le lot figé, avec la vérification $C_1$--$C_4$
appliquée à tous de la même manière.

\begin{thebibliography}{9}
\bibitem{dinkelbach1967}
W.~Dinkelbach.
\newblock On nonlinear fractional programming.
\newblock \emph{Management Science}, 13(7):492--498, 1967.

\bibitem{zerdani2011}
O.~Zerdani et M.~Moulaï.
\newblock Optimization over an integer efficient set of a multiple objective
  linear fractional problem.
\newblock \emph{Applied Mathematical Sciences}, 5(50):2451--2466, 2011.

\bibitem{drici2018}
W.~Drici \emph{et al.}
\newblock Optimisation sur l'ensemble efficace entier d'un programme
  fractionnaire multi-objectif, 2018.
\newblock \emph{[Référence citée d'après le plan de thèse ; coordonnées
  bibliographiques complètes à vérifier sur l'original.]}

\bibitem{sylva2004}
J.~Sylva et A.~Crema.
\newblock A method for finding the set of non-dominated vectors for multiple
  objective integer linear programs.
\newblock \emph{European Journal of Operational Research}, 158(1):46--55, 2004.
\end{thebibliography}

\end{document}
